\section{Lecture 6 (September 17th)}
\begin{ex}
($U(1)$ symmetry in a complex scalar field theory) Consider the Lagrangian 
\[\mathcal{L}=-\partial ^{\mu }\phi \partial _{\mu }\phi ^{*}-\mathcal{V}(|\phi |^2)\]
The $U(1)$ symmetry is given as
\[\phi '(x)=e^{i\alpha }\phi (x)=\phi (x)+\underbrace{i\alpha \phi (x)}_{\delta \phi (x)}+O(\alpha ^2)\]
Invariance of the action is given by $\delta \mathcal{L}=\partial _{\mu }f^{\mu }=0\rightarrow f^{\mu }=0$. The Noether current is
\[\begin{align*}
j^{\mu }=\dfrac{\partial \mathcal{L}}{\partial (\partial _{\mu }\phi )}\,\delta \phi + \dfrac{\partial \mathcal{L}}{\partial (\partial _{\mu }\phi ^{*})}\,\delta \phi ^{*}-f^{\mu }=(-\partial ^{\mu }\phi ^{*})(i\alpha \phi )+(-\partial ^{\mu }\phi )(-i\alpha \phi ^{*})\\
=&\alpha (-i\phi \partial ^{\mu }\phi ^{*}+i\phi ^{*}\partial ^{\mu }\phi )
\end{align*}\]
which implies that $\partial _{\mu }j^{\mu }=0$ on-shell. Then, Noether charge would be, physically, the charge carried by particles described by $\phi $.
\end{ex}
\vspace{2ex}
\begin{ex}
($U(1)$ symmetries in Dirac theory) Consider the Lagragian
\[\mathcal{L}=\bar{\psi  }(i\cancel{\partial }-m)\psi \]
The $U(1)$ symmetry is given as
\[\psi '(x)=e^{i\alpha }\psi (x)=\psi (x)+\underbrace{i\alpha \psi (x)}_{\delta \psi (x)}+O(\alpha ^2)\]
Again, $\delta \mathcal{L}=0$ implies that $f^{\mu }=0$ and the Noether current is given by
\[j^{\mu }=\dfrac{\partial \mathcal{L}}{\partial (\partial _{\mu }\psi )}\,\delta \psi +\dfrac{\partial \mathcal{L}}{\partial (\partial _{\mu }\psi ^{\dagger})}\,\delta \psi^{\dagger}-f^{\mu }=i\bar{\psi}\gamma ^{\mu }(i\alpha \psi )  \]
and $\partial _{\mu }j^{\mu }=0$. The rather annoying sign can be eliminated via renormalisation giving
\[j^{\mu }=\bar{\psi }\gamma ^{\mu }\psi \]
\end{ex}
\vspace{2ex}
\begin{ex}
(Chiral transformation) Another $U(1)$ symmetry arises for the massless case $m=0$.
\[\psi '(x)=e^{i\alpha \gamma ^{5}}\psi (x)=\psi (x)+i\alpha \gamma ^{5}\psi (x)+O(\alpha^2)\]
The Nother current is given as, through renormalising $\delta \psi =-i\gamma ^{5}\psi $
\[j^{\mu }_{5}=\bar{\psi }\gamma ^{\mu }\gamma ^{5}\psi \]
We call $j^{\mu }$ and $j^{\mu}_{5}$ are called vectors and axial-vector current respectively. For $\psi $ coupled to an electromagnetic field (EM) $A_{\mu }$, $j^{\mu }$ is called the electric current density. For a massless $\psi $ coupled with a EM field,
\[\begin{cases}
j^{\mu }_{L}=\dfrac{1}{2}(j^{\mu }-j_{5}^{\mu })=\bar{\psi }\gamma ^{\mu }\dfrac{1-\gamma ^{5}}{2}\psi =\bar{\psi }_{L}\gamma ^{\mu }\psi _{L}\\
j^{\mu }_{R}=\dfrac{1}{2}(j^{\mu }+j^{\mu }_{5})=\bar{\psi }\gamma ^{\mu }\dfrac{1+\gamma ^{5}}{2}\psi =\bar{\psi }_{R}\gamma ^{\mu }\psi _{R}
\end{cases}\]
and these denote the electric current densities for left-handed and right-handed Weyl spinors. 
\end{ex}
\vspace{2ex}
\begin{ex}
(Translation symmetry in a real scalar field) Consider the Lagrangian
\[\mathcal{L}=-\dfrac{1}{2}\partial ^{\mu }\phi \partial _{\mu }\phi -\mathcal{V}(\phi )\]
We need to see what $\delta \phi $ is under $x\rightarrow x'=x+a$. By definition, 
\[\phi '(x')=\phi (x)=\phi (x'-a)=\phi (x')-\underbrace{a^{\mu }\partial _{\mu }\phi (x')}_{\delta \phi (x')}+O(a^2)\]
Now what is $f^{\mu }$? We can easily expect $\delta \mathcal{L}=\mathcal{L}'-\mathcal{L}=-a^{\mu }\partial _{\mu }\mathcal{L}$. Therefore, $f^{\mu }=-a^{\mu }\mathcal{L}$. Our Noether current is
\[j^{\mu }=\dfrac{\partial \mathcal{L}}{\partial (\partial _{\mu }\phi )}\,\delta \phi -f^{\mu }=a_{\nu }(\partial ^{\mu }\partial ^{\nu }\phi +\eta ^{\mu \nu }\mathcal{L}) \]
The rank 2 tensor in the second term is called the stress tensor $T^{\mu \nu }$. We find that
\[\partial _{\mu }j^{\mu }=0=\partial _{\mu }(a_{\nu }T^{\mu \nu })\]
for all $a^{\mu }$ and 
\[\partial _{\mu }T^{\mu \nu }=0\]
It is natural to think that the conserved quantity has four degrees of freedom as we have translated with four degrees of freedom. The Noether charges exactly the 4-momentum!
\[P^{\nu }=\int d^3\vec{x}\,T^{0\nu }\]
We can then check that
\[P^{0}=H=\int d^3\vec{x}(\pi \dot{\phi }-\mathcal{L})\]
Simply beautiful!
\end{ex}
\vspace{2ex}
\begin{ex}
(Lorentz symmetry in a real scalar field theory) Consider the infinitesimal transformation $x'=\mathrm{\Lambda }x=(1+\omega )x$. Componentwise, we can write
\[x'^{\mu }=\mathrm{\Lambda }^{\mu }{}_{\nu }x^{\nu }=(\delta ^{\mu }{}_{\nu }+\omega ^{\mu }{}_{\nu })x^{\nu }\]
Using the Taylor expansion again,
\[\phi'(x')=\phi (x)=\phi (\mathrm{\Lambda }^{-1}x')=\phi (x')-\omega ^{\mu }{}_{\nu }x'^{\nu }\partial _{\mu }\phi (x')+O(\omega ^2)\]
the trick is that $(1+\omega )(1-\omega )=1+O(\omega ^2)$. What is $f^{\mu }$?
\[\delta \mathcal{L}=\mathcal{L}'-\mathcal{L}=-\omega ^{\mu }{}_{\nu }x^{\nu }\partial _{\mu }\mathcal{L}=\partial _{\mu }(-\omega ^{\mu }{}_{\nu }x^{\nu }\mathcal{L})=\partial _{\mu }f^{\mu }\]
Thus, $f^{\mu }=-\omega ^{\mu }{}_{\nu }x^{\nu }\mathcal{L}$. The Noether current is correspondingly
\[\begin{align*}
j^{\mu }=&\dfrac{\partial \mathcal{L}}{\partial (\partial _{\mu }\phi )}\,\delta \phi -f^{\mu }\\
=&\omega ^{\nu \rho }x_{\rho }\partial _{\nu }\phi \partial ^{\mu }\phi +\omega ^{\nu \rho }x_{\rho }\delta ^{\mu }{}_{\nu }\mathcal{L}\\
=&\omega _{\nu \rho }x^{\rho }\underbrace{(\partial ^{\nu }\phi \partial ^{\mu }\phi +\eta ^{\mu \nu }\mathcal{L})}_{T^{\mu \nu }}\\
=&\dfrac{1}{2}\omega _{\nu \rho }(x^{\rho }T^{\mu \nu }-x^{\nu }T^{\mu \rho })
\end{align*}\]
we call the second term $\mathcal{J}^{\mu \nu \rho }$. Concluding,
\[\partial _{\mu }j^{\mu }=0\]
for all $\omega _{\nu \rho }$ implies that
\[\partial _{\mu }\mathcal{J}^{\mu \nu \rho }=0\]
Thus, the Noether charges are generalised (with boosts) angular momenta!
\[J^{\mu \nu }=\int d^3\vec{x}\,\mathcal{J}^{0\mu \nu }\]
for example, $\mu \nu =12$ shows
\[\int d^3\vec{x}\,(x^2T^{01}-x^{1}T^{02})\]
\end{ex}
\vspace{2ex}
\begin{rmk}
There are also discrete symmetries that we learn from here.
\end{rmk}
\vspace{2ex}
\begin{defi}
(Space inversion or parity) The transformation is given by
\[x=(t,\vec{x})=(t',\vec{x}')=(t,-\vec{x})\]
The parity on fields give
\[\phi '(x')=\phi (x)\qquad\&\qquad \phi '(x')=\gamma ^{0}\psi (x)\]
lastly,
\[A'^{\mu }(x')=\begin{cases}
A^{\mu }(x)&(\mu =0)\\
-A^{\mu }(x)&(\mu =1,2,3)
\end{cases}\]
\end{defi}
\vspace{2ex}
\begin{ex}
(Parity on $U(1)$ currents in Dirac theory) For vectors,
\[j^{\mu }=\bar{\psi }\gamma ^{\mu }\psi \rightarrow (\gamma ^{0}\psi )^{\dagger}\gamma ^{0}\gamma ^{\mu }(\gamma ^{0}\psi )=\psi ^{\dagger}\gamma ^{\mu }\gamma ^{0}\psi \]
and
\[\begin{cases}
 \bar{\psi }\gamma ^{0}\psi &(\mu =0)\\
-\bar{\psi }\gamma ^{i}\psi &(\mu =i)
\end{cases}\]
while for axial-vectors,
\[j^{\mu }_{5}=\bar{\psi }\gamma ^{\mu }\gamma _{5}\psi \rightarrow (\gamma ^{0}\psi )^{\dagger}\gamma ^{0}\gamma ^{\mu }\gamma _{5}(\gamma ^{0}\psi )=\psi ^{\dagger}\gamma ^{\mu }\gamma _{5}\gamma ^{0}\psi \]
and
\[\begin{cases}
-\bar{\psi }\gamma ^{0}\gamma _{5}\psi &(\mu =0)\\
\bar{\psi }\gamma ^{\mu }\gamma _{5}\psi &(\mu =i)
\end{cases}\]
\end{ex}
\vspace{2ex}
\begin{defi}
(Charge conjugation) This transformation simply swaps particles with anti-particles and has nothing to do with position. Charge conjugation on fields result in
\[\phi '(x)=\phi (x)\& \phi '(x)=C\gamma ^{0}\psi ^{*}(x)\& A_{\mu }'(x)=-A_{\mu }(x)\]
here,
\[C\gamma ^{\mu }C^{-1}=-(\gamma ^{\mu })^{T}\& C^{-1}=C^{\dagger}=C^{T}=-C\]
and a typical choice is
\[C=i\gamma ^2\gamma ^{0}\]
\end{defi}
\vspace{2ex}
\begin{defi}
(Time reversal) The time reversal transformation of $x$ is given as $x=(t,\vec{x})\rightarrow (t',\vec{x}')=(-t,\vec{x})$. The time reversal on fields give
\[\phi '(x')=\phi (x)^{*}\qquad\&\qquad \psi '(x')=T\psi (x)^{*}\]
and finally
\[A'^{\mu }(x')=\begin{cases}
A^{\mu }(x)&(\mu =0)\\
-A^{\mu }(x)&(\mu =i)
\end{cases}\]
Here, $T$ is the time-reversal matrix, satisfying 
$T^{-1}\gamma ^{\mu }T=(\gamma ^{\mu })^{T}\qquad \&\qquad T=T^{-1}=T^{\dagger}=-T^{*}=-T^{T}$
again, a typical choice is
\[T=i\gamma ^{1}\gamma ^{3}=-\gamma ^{5}C\]
\end{defi}
\vspace{2ex}
 
